\documentclass[11pt]{article}

%====================
% Basic setup
%====================
\usepackage[left=0.7in, right=0.7in, top=0.7in, bottom=0.7in]{geometry}
\usepackage{hyperref}
\usepackage{xcolor}
\usepackage{enumitem}
\usepackage{titlesec}
\usepackage{setspace}

\hypersetup{
    colorlinks=true,
    urlcolor=black
}

\setlength{\parindent}{0pt}
\setlength{\parskip}{4pt}
\pagenumbering{gobble}






%====================
% Section formatting
%====================
\titleformat{\section}
  {\large\bfseries}
  {}
  {0pt}
  {}
  [\vspace{2pt}\titlerule]

%====================
% Document
%====================
\begin{document}





%====================
% Header
%====================
{\Huge \textbf{Angel Fernando Bórquez Guerrero}} \\
\vspace{-0.4cm}
\par \href{https://www.linkedin.com/in/fernandoborquez/}{LinkedIn} \quad | \quad
\href{https://github.com/FernyBoy}{GitHub} \quad | \quad
\href{mailto:fernandoborquez215@icloud.com}{fernandoborquez215@icloud.com} \quad | \quad
+52 633 100 5991 \\
Hermosillo, Sonora.





%====================
% Educación
%====================
\section*{Educación}
\textbf{Universidad de Sonora}
Licenciatura en Ciencias de la Computación \hfill \textbf{August 2023 -- Present}



%====================
% Experiencia Profesional
%====================
\section*{Experiencia Profesional}
\textbf{TE Connectivity} \hfill Supervisor: M.Sc. Martín Vega \\
Becario de Ingeniería de IA \hfill \textbf{Junio 2026 -- Septiembre 2026}
\vspace{-0.3cm}
\begin{itemize}[leftmargin=*]
    \item Diseñé y desarrollé un agente de IA para el análisis de alertas industriales, implementando su orquestación, gestión de estados, integración de herramientas y flujos de ejecución para apoyar el análisis y la toma de decisiones automatizados.

    \vspace{-0.3cm}
    \item Rediseñé la arquitectura de orquestación del agente utilizando LangGraph, reemplazando la implementación anterior del flujo de trabajo por una arquitectura basada en grafos que mejoró la gestión de estados, la modularidad y la eficiencia de ejecución.

    \vspace{-0.3cm}
    \item Optimicé la arquitectura y el flujo de ejecución del agente, reduciendo los tiempos de respuesta de hasta tres minutos a respuestas prácticamente instantáneas.

    \vspace{-0.3cm}
    \item Diseñé e implementé un servidor Linux compartido para becarios, configurando el sistema y el acceso de usuarios para proporcionar un entorno centralizado para desarrollar y alojar proyectos de investigación e ingeniería.
\end{itemize}


\vspace{0.2cm}
\textbf{ACARUS — Área de supercómputo} \hfill Supervisor: Dr. María del Carmen Heras Sánchez\\
Becario Profesional \hfill \textbf{Septiembre 2025 -- Presente}
\vspace{-0.3cm}
\begin{itemize}[leftmargin=*]
    \item Realicé una estancia profesional en el área de Supercómputo de la Universidad de Sonora.
    
    \vspace{-0.3cm}
    \item Formé parte de un equipo de desarrollo encargado de construir un foro especializado para los usuarios de \textit{Yuca}, el sistema HPC de la universidad.
    
    \vspace{-0.3cm}
    \item Implementé funcionalidades como publicaciones de usuarios, mensajería directa y flujos de solicitudes de servicio.
    
\end{itemize}


%====================
% Experiencia en Investigación
%====================
\section*{Experiencia en Investigación}
\textbf{University of Guadalajara} \hfill Guadalajara, México \\
\textit{Becario de Investigación} (Asesor: Dr. Rafael Morales Gamboa) \hfill \textbf{Junio 2025 -- Julio 2025}
\vspace{-0.3cm}
\begin{itemize}
    \item Validé y generalicé un modelo de Memoria Asociativa Entrópica (EAM) en conjuntos de datos de referencia de alta dimensionalidad, culminando en una publicación arbitrada en \textit{IEEE ENC 2025}.

    \vspace{-0.3cm}
    \item Diseñé un marco experimental de detección de novedades para evaluar rigurosamente el desempeño de rechazo de clases fuera de distribución (OOD) de la arquitectura EAM bajo diferentes condiciones operativas.

    \vspace{-0.3cm}
    \item Integré redes perceptuales profundas (autoencoders convolucionales y clasificadores) para proyectar espacios de imágenes sin procesar en representaciones vectoriales latentes acotadas compatibles con topologías de memoria asociativa.

    \vspace{-0.3cm}
    \item Formulé protocolos de control cuantitativo utilizando métricas especializadas de ``no respuesta'' para modelar la saturación parcial de la memoria y caracterizar dinámicas sistemáticas de rechazo de falsos positivos.

    \vspace{-0.3cm}
    \item Desarrollé un marco de simulación reproducible de extremo a extremo, modularizando distribuciones de densidad de clases, proporciones de balance y esquemas de evaluación de recuperación multiclase.
\end{itemize}


%====================
% Publicaciones
%====================
\section*{Publicaciones}
\begin{itemize}
    \item \textbf{A. F. Bórquez Guerrero}, R. Morales Gamboa, et al. ``\textit{Large-Scale Validation and Analysis of the Rejection Capacity in Entropic Associative Memory}.'' In \textbf{Proceedings of the 2025 Mexican International Conference on Computer Science (ENC) \& IEEE ICEV}, Orizaba, Veracruz, México, 2025. IEEE. \\
    \footnotesize\textit{Indexing: Online ISSN: 2332-5712 $\vert$ ISBN: 979-8-3315-9026-0/25 $\vert$ IEEE Xplore.}
\end{itemize}


%====================
% Competitions & Projects
%====================
\section*{Competencias \& Proyectos}
\textbf{TE AI CUP 2026 -- 1.er Lugar Internacional} \hfill Proyecto: Industrial Alert Agent\\
Líder de Proyecto \hfill \textbf{Enero 2026 - Junio 2026}
\vspace{-0.3cm}
\begin{itemize}[leftmargin=*]
    \item Lideré un equipo interdisciplinario de 6 integrantes en una competencia internacional de IA, dirigiendo el desarrollo técnico de un sistema basado en IA para el análisis de alertas industriales que obtuvo el 1.er lugar. 

    \vspace{-0.3cm} 
    \item Desarrollé un agente de IA que integra Modelos de Lenguaje de Gran Escala (LLMs) para el análisis automatizado de alertas y apoyo a la toma de decisiones.

    \vspace{-0.3cm} 
    \item Desarrollé un pipeline de datos y aprendizaje automático de extremo a extremo que incluyó limpieza de datos, preprocesamiento de texto, generación de embeddings, clustering, validación de datos y modelado predictivo. 

    \vspace{-0.3cm}
    \item Desarrollé un sistema de modelado predictivo para estimar la probabilidad de futuras alertas industriales, proporcionando un componente proactivo que complementó el análisis de patrones históricos y permitió identificar posibles problemas con mayor anticipación.

    \vspace{-0.3cm} 
    \item Diseñé la arquitectura del proyecto alrededor de estructuras de software limpias, estandarizadas y escalables, estableciendo convenciones para módulos, herramientas, configuración, gestión de estados e integración de componentes. 
\end{itemize}



\textbf{HSIL Hackathon 2026 -- 3.er Lugar} \hfill Proyecto: ShiftGuardian\\
Diseñador de Aplicaciones \& Arquitecto de Agentes de IA \hfill \textbf{Abril 2026}
\vspace{-0.3cm}
\begin{itemize}[leftmargin=*]
    \item Diseñé la interfaz de la aplicación y la arquitectura del agente de IA para un sistema de apoyo a la toma de decisiones clínicas orientado a la transferencia de pacientes entre enfermeros y su priorización.
    
    \vspace{-0.3cm}
    \item Integré un concepto de flujo de trabajo basado en RAG y LLM para resumir reportes clínicos, identificar tareas críticas y reducir la carga cognitiva durante los cambios de turno.
\end{itemize}



\textbf{MICAI 2025 MeDA Challenge -- 1.er Lugar Nacional} \hfill \textbf{Octubre 2025}
\vspace{-0.3cm}
\begin{itemize}[leftmargin=*]
    \item Contribuí al 1.er lugar en una competencia académica nacional enfocada en Adaptación de Dominio Médico mediante Aprendizaje Auto-Supervisado (SSL), lo que resultó en una invitación para publicar en una edición especial de la revista \textit{Health Information Science and Systems}.
    
    \vspace{-0.3cm}
    \item Estructuré y estandaricé el conjunto de datos y los recursos de entrenamiento para establecer un pipeline reproducible que facilitara el desarrollo y la experimentación con modelos.
    
    \vspace{-0.3cm}
    \item Desarrollé recursos visuales y de presentación utilizados para comunicar el proyecto durante las etapas de competencia y divulgación.
\end{itemize}



\textbf{TE AI CUP 2025 -- 3.er Lugar Internacional} \hfill Proyecto: AI Generated Machine Downtime Insights\\
Desarrollador UI \hfill \textbf{Enero -- Mayo 2025}
\vspace{-0.3cm}
\begin{itemize}[leftmargin=*]
    \item Diseñé y desarrollé la UI/UX de una plataforma basada en IA para el análisis de tiempos de inactividad industrial, traduciendo datos complejos de series temporales de máquinas y resultados de aprendizaje no supervisado en información visual clara y accionable.
    
    \vspace{-0.3cm}
    \item Desarrollé una interfaz accesible y centrada en el usuario para personas no técnicas, simplificando la exploración e interpretación de resultados generados por IA sin requerir conocimientos previos de aprendizaje automático.
    
\end{itemize}



\end{document}
